\chapter{Arquitectura e implementación de la solución perimetral}
\label{chapter:chapter05}

Este capítulo toma la selección metodológica del Capítulo~\ref{chapter:chapter04} y la convierte en una arquitectura perimetral concreta para la Finca La Esmeralda. Primero se justifica la alternativa adoptada; luego se describe la base física de captura, el servicio residente en el borde y el motor que organiza la evidencia consolidada. La validación de desempeño se reserva para el Capítulo~\ref{chapter:chapter06}.

\section{Análisis de soluciones y selección final}
La solución perimetral parte de la selección metodológica del Capítulo~\ref{chapter:chapter04}. A partir de ese punto, esta sección sintetiza las alternativas consideradas y los criterios que llevaron a la configuración adoptada.

\subsection{Evaluación técnica de alternativas de solución}
Para atender la pérdida de datos y conservar la relación entre telemetría operativa (\acs{RS-232}) y biótica (\acs{UDP}), se analizaron cuatro arquitecturas bajo los criterios de continuidad, no invasividad, soberanía local de datos y compatibilidad con \textit{SenseHub Dairy}.

\begin{enumerate}
    \item \textbf{Propuesta 1: Redireccionamiento lógico mediante agentes de programa (virtualización):}

    Esta opción evita intercepción física y concentra la captura dentro de la computadora de la lechería, mediante puertos virtuales para el flujo serial y bibliotecas de captura sobre la interfaz Ethernet original. Aunque reduce cableado, exige privilegios administrativos, modo promiscuo e instalación de controladores; por tanto, contradice la condición de caja negra del sistema comercial.

    \item \textbf{Propuesta 2: Duplicación mediante caja derivadora comercial de arquitectura cerrada:}

    Esta alternativa replica la señal con equipos terminados, como derivadores \textit{RS-232} y concentradores Ethernet convencionales. Simplifica la instalación, pero su diseño cerrado impide ajustar impedancias o incorporar aislamiento galvánico; por ello mantiene riesgo de lazos de tierra y dominios de colisión.

    \item \textbf{Propuesta 3: Acondicionamiento activo mediante MAX3232, cable en Y y derivación pasiva de red:}

    Esta propuesta introduce una derivación física en Y y acondicionamiento con MAX3232 para trasladar niveles eléctricos hacia el nodo de captura. Mejora la lectura serial frente a una derivación puramente pasiva, pero no incorpora aislamiento industrial y el acoplamiento pasivo de red puede introducir desadaptaciones de impedancia.

    \item \textbf{Propuesta 4: Interfaz industrial aislada (Waveshare), cable en Y y conmutación de Capa 2:}

    Esta arquitectura observa el bus serial mediante cable en Y e interfaz industrial Waveshare con aislamiento galvánico, mientras el tráfico de collares se separa con conmutación de Capa 2. Ofrece soberanía sobre ambos dominios de evidencia sin intervenir la computadora de la lechería ni asumir control sobre la lógica de \textit{SenseHub Dairy}.
\end{enumerate}

\subsection{Selección de la arquitectura y criterios de diseño}
Las alternativas se confrontaron mediante una matriz de Pugh con cinco criterios ponderados ($W_i$) y una escala de 1 a 5, tomando la \textbf{Propuesta 2} como referencia convencional de mercado.

\begin{table}[H]
\centering
\caption{Matriz de selección de alternativas tecnológicas (Matriz de Pugh).}
\label{tab:matriz_pugh}
\footnotesize
\setlength{\tabcolsep}{4pt}
\renewcommand{\arraystretch}{1.25}
\begin{tabularx}{\linewidth}{|>{\raggedright\arraybackslash}X|>{\centering\arraybackslash}p{1.35cm}|>{\centering\arraybackslash}p{0.85cm}|>{\centering\arraybackslash}p{1.15cm}|>{\centering\arraybackslash}p{0.85cm}|>{\centering\arraybackslash}p{0.85cm}|}
\hline
\rowcolor{gray!20} \textbf{Criterios de evaluación} & \textbf{Peso ($W_i$)} & \textbf{P1} & \textbf{P2 (ref.)} & \textbf{P3} & \textbf{P4} \\ \hline
C1. Aislamiento galvánico y protección \ac{EMC} & 25\% & 1 & 1 & 2 & 5 \\ \hline
C2. Gestión de red (\acs{UDP} sin colisiones) & 25\% & 2 & 1 & 2 & 5 \\ \hline
C3. Soberanía e independencia administrativa & 20\% & 1 & 5 & 5 & 5 \\ \hline
C4. No invasividad (garantía Allflex) & 15\% & 1 & 3 & 4 & 5 \\ \hline
C5. Determinismo y latencia (\acs{RS-232}) & 15\% & 1 & 3 & 5 & 5 \\ \hline \hline
\textbf{Puntaje total ponderado} & \textbf{100\%} & \textbf{1{,}25} & \textbf{2{,}40} & \textbf{3{,}35} & \textbf{5{,}00} \\ \hline
\rowcolor{ThesisGreen!8} \textbf{Orden de clasificación} & -- & \textbf{4} & \textbf{3} & \textbf{2} & \textbf{1} \\ \hline
\end{tabularx}
\end{table}

La matriz favorece la \textbf{Propuesta 4} como alternativa para esta intervención. La decisión se sostiene en cuatro criterios: menor riesgo eléctrico sobre el ecosistema original, mayor control del dominio de red, independencia frente a restricciones administrativas y mejores condiciones para capturar en presencia de ruido electromagnético. A partir de esa selección, la escucha pasiva del canal serial se definió en \textbf{19\,200\,bps}, configuración que reduce la ventana temporal de exposición al ruido, disminuye el tiempo de ocupación por trama y aumenta la capacidad nominal de transporte. Las secciones siguientes desarrollan el alcance funcional y la implementación física y lógica de la solución seleccionada.

\section{Base física de la solución}

\subsection{Alcance funcional de la pasarela perimetral}
Tal como se estableció en la sección~\ref{sec:diseno-integracion-perimetral}, la pasarela perimetral concreta cuatro decisiones de diseño: observar sin invadir, aislar eléctricamente, separar el tráfico útil y conservar evidencia cuando la publicación no es posible. Esta sección distingue esa infraestructura de captura del motor de análisis, que opera después sobre evidencia ya consolidada.

La pasarela perimetral se implementa sobre un nodo de borde de intervención mínima. Su función es \textbf{desacoplar} el flujo telemétrico del ecosistema \textit{SenseHub Dairy} de perturbaciones del entorno, como el ruido electromagnético, la variación temporal de la latencia y el tráfico administrativo en la \ac{LAN}, sin modificar la lógica original del sistema ni comprometer la integridad del bus de control. En la práctica, actúa como un sistema de \textbf{adquisición, selección determinista de tramas relevantes, validación de integridad y sincronización controlada}.

Su operación se organiza en cuatro etapas funcionales, sintetizadas en la Tabla~\ref{tab:etapas_funcionales_pasarela}. La tabla permite leer la pasarela como una cadena de decisiones: captura no invasiva, separación de dominios, normalización controlada y persistencia ante fallos.

\begin{table}[H]
\centering
\caption{Etapas funcionales de la pasarela perimetral.}
\label{tab:etapas_funcionales_pasarela}
\footnotesize
\setlength{\tabcolsep}{4pt}
\renewcommand{\arraystretch}{1.08}
\begin{tabularx}{\linewidth}{|>{\raggedright\arraybackslash}p{3.15cm}|>{\raggedright\arraybackslash}X|>{\raggedright\arraybackslash}X|}
\hline
\rowcolor{gray!20} \TableHeaderFirst{Etapa} & \TableHeader{Decisión de diseño} & \TableHeader{Función en la pasarela} \\ \hline
Adquisición no invasiva & Derivación \ac{RS-232} en modo escucha, cable en Y e interfaz aislada. & Capturar la señal sin cargar el bus original y reducir la exposición a lazos de tierra o acoplamientos de modo común. \\ \hline
Separación por dominio & Distinción temprana entre ordeño completo, telemetría de collar y evidencia de red. & Evitar mezclar señales con significados distintos antes de cualquier correlación o transporte. \\ \hline
Normalización y política de borde & Metadatos temporales, mensajes \acs{JSON} y lista de permitidos por IP, puerto y protocolo. & Convertir la evidencia aceptada en registros trazables y retener lotes que no satisfacen la política. \\ \hline
Persistencia y transporte & Tópicos \ac{MQTT}, espejo local y cola \textit{JSONL} de almacenamiento y reenvío. & Preservar continuidad cuando el destino no está disponible o la publicación debe diferirse. \\ \hline
\end{tabularx}
\end{table}

En conjunto, la pasarela actúa como un \textbf{filtro determinista de borde} y un \textbf{amortiguador de continuidad}: separa dominios, normaliza evidencia, aplica reglas simples de seguridad y mantiene persistencia local para desacoplar la captura de la variabilidad del entorno.

Las secciones siguientes desarrollan la implementación física, lógica y operativa de estas funciones. Más adelante, la sección del motor de análisis se reserva para el tratamiento forense y de correlación de la evidencia ya capturada por la pasarela.

\subsection[Justificación técnica de la tasa de transferencia]{Justificación técnica de la tasa de transferencia\\(\texorpdfstring{19\,200\,bps}{19200 bps})}
La tasa de transferencia de la pasarela se definió como parte del diseño, no como un ajuste aislado. La escucha pasiva se sincronizó a \textbf{19\,200\,bps} para equilibrar integridad y eficiencia por dos razones principales:

\begin{itemize}
    \item \textbf{Reducción de la ventana de exposición al ruido:} al duplicar la tasa de baudios, el tiempo de bit ($T_{bit}$) se reduce de $104{,}17\,\mu\text{s}$ a \textbf{$52{,}08\,\mu\text{s}$}. Esta compresión temporal disminuye estadísticamente la probabilidad de que un transitorio electromagnético (pico de \acs{EMI}) coincida con la fase crítica de muestreo del bit.
    \item \textbf{Optimización de la eficiencia del canal ($\eta$):} una mayor velocidad reduce el tiempo de permanencia de cada trama en el bus y deja más margen para los procesos administrativos del fabricante. La mejora relativa indicada en la Tabla~\ref{tab:comparativa_baudios_detallada} se usa como criterio de diseño, no como resultado experimental.
\end{itemize}

\begin{table}[H]
\centering
\caption{Comparativa técnica de temporización: 9600 vs. 19\,200 bps.}
\label{tab:comparativa_baudios_detallada}
\footnotesize
\renewcommand{\arraystretch}{1.3}
\begin{tabular}{|l|c|c|c|}
\hline
\rowcolor{gray!20} \TableHeaderFirst{Parámetro de diseño} & \TableHeader{9600 bps} & \TableHeader{19\,200 bps} & \TableHeader{Mejora relativa} \\ \hline
Tiempo de bit ($T_{bit}$) & $104{,}17\,\mu\text{s}$ & \textbf{$52{,}08\,\mu\text{s}$} & -50\% de exposición \\ \hline
Tiempo de trama (10 bits) & $1{,}04\,\text{ms}$ & \textbf{$0{,}52\,\text{ms}$} & Mayor disponibilidad \\ \hline
Capacidad nominal relativa ($\varepsilon$) & $\approx$ 0{,}28\% & \textbf{$\approx$ 0{,}55\%} & +96\% de rendimiento \\ \hline
\end{tabular}
\end{table}

Los valores de temporización presentados en la Tabla~\ref{tab:comparativa_baudios_detallada} se calculan sobre una trama asíncrona de 10 bits (1 de inicio, 8 de datos y 1 de parada). En esta etapa se emplean como justificación de diseño de la pasarela; la validación experimental de su efecto sobre eficiencia, pérdida y estabilidad del sistema se expone en el Capítulo~\ref{chapter:chapter06}.

\subsection{Arquitectura física de la pasarela perimetral (Niveles 0 y 1)}
Para formalizar el diseño, se utilizó la metodología de modelado funcional \textbf{IDEF0}. En la Figura \ref{fig:nivel0} se establece el \textbf{Diagrama de contexto (Nivel 0)}, para definir las fronteras operativas del sistema bajo el estándar \textbf{ICOM} (Entradas, Controles, Mecanismos y Salidas).

\begin{figure}[H]
    \centering
    \includegraphics[width=0.86\textwidth]{Figuras/nivel_0.png} 
    \caption{Diagrama de contexto Nivel 0: Pasarela Perimetral de Datos Bióticos (elaboración propia).}
    \label{fig:nivel0}
\end{figure}

\subsubsection{Nivel 0 bajo el estándar ICOM}
El bloque \textbf{A0} modela la pasarela como una unidad de procesamiento integral. Su funcionamiento se define mediante los cuatro componentes del estándar:
\begin{itemize}
    \item \textbf{Entradas:} trama serial inalterada proveniente del pin 2 (RX) del SenseHub y tráfico Ethernet UDP del puerto 6001 de los collares, extraídos mediante el conmutador de Capa 2.
    \item \textbf{Controles:} operaci\'{o}n bajo las condiciones de red observadas en la l\'{i}nea base, con un desfase temporal m\'{a}ximo de referencia de $71{,}3\,\text{ms}$, uso previsto de \acs{MQTT} como mecanismo de mensajer\'{i}a local dentro del per\'{i}metro de la finca y soporte de \acs{TLS} 1.3 para proteger la telemetr\'{i}a en tr\'{a}nsito hacia el intermediario local de borde.
    \item \textbf{Mecanismos:} módulo de aislamiento galvánico (\textit{Waveshare}) con barrera superior a $3{,}0\,\text{kVrms}$, conmutador de Capa 2 (\textit{NETGEAR}), procesador central (Raspberry Pi 5), sistema de alimentación ininterrumpida (\acs{UPS}) APC-1500 para resiliencia energética ante fallos del suministro eléctrico y alimentación de 120 V AC.
    \item \textbf{Salidas:} telemetr\'{i}a estructurada en \acs{JSON} hacia el intermediario local de borde y registros locales que conservan el contexto de captura para revisión posterior.
\end{itemize}

Al expandir el sistema, el \textbf{Diagrama Funcional de Nivel 1} (Figura \ref{fig:nivel1}) descompone la operación en tres etapas secuenciales diferenciadas:

\begin{figure}[H]
    \centering
    \includegraphics[width=0.86\textwidth]{Figuras/nivel1.png} 
    \caption{Diagrama Funcional Nivel 1: Descomposición funcional de la pasarela (elaboración propia).}
    \label{fig:nivel1}
\end{figure}

\subsubsection{Subprocesos del Nivel 1}
A diferencia de un puente de comunicación convencional, la pasarela ejecuta tareas de protección y procesamiento selectivo en cascada:
\begin{itemize}
    \item \textbf{A1 (Interceptar y Aislar):} representa la frontera física de captura en los dos dominios de telemetría. En el canal serial, utiliza un cable en Y para extraer la señal \acs{RS-232} sin carga resistiva sobre el bus original; el convertidor industrial \textbf{Waveshare} traslada los niveles eléctricos ($\pm$12\,V) a niveles lógicos (3{,}3\,V), proporcionando desacoplo magnético y contribuyendo a proteger al procesador frente a transitorios mediante una barrera de aislamiento de grado industrial. En el dominio Ethernet, incorpora un \textbf{conmutador de Capa 2} que mantiene el reenvío por direcciones \acs{MAC}, de modo que el tráfico \acs{UDP} de los collares conserva una ruta operativa segmentada sin insertar carga ni alterar el sistema \textit{SenseHub}.
    \item \textbf{A2 (Procesar e Inspeccionar):} Concentra el procesamiento selectivo. Recibe la señal serial digital aislada y el flujo Ethernet segmentado proveniente del conmutador de Capa 2, y aplica el módulo de inspección profunda (\acs{DPI}). Su función es actuar como un filtro activo que distingue durante el ciclo de captura entre ráfagas bióticas útiles, eventos seriales de ordeño y tráfico administrativo no prioritario, optimizando los recursos antes del transporte.
    \item \textbf{A3 (Encapsular y Sincronizar):} gestiona la capa de mensajer\'{i}a y sincronizaci\'{o}n. Recibe los datos validados y consolidados para transformarlos en objetos \acs{JSON}, y habilita su publicaci\'{o}n as\'{i}ncrona hacia el intermediario local de borde (Mosquitto sobre Raspberry Pi) o hacia un espejo local en disco, con soporte de \acs{TLS} 1.3 para protecci\'{o}n en tr\'{a}nsito dentro del per\'{i}metro.
\end{itemize}

Como integración física de estos subprocesos, la pasarela completa se aloja en un gabinete \textbf{IP65}, lo que limita la exposición de los componentes de captura, procesamiento y transporte a humedad y salpicaduras propias de una finca.

\section{Componente lógico de la solución}

\subsection{Arquitectura lógica de la pasarela perimetral (Nivel 2)}
El bloque \textbf{A2} representa la lógica de ejecución en el borde. La Figura \ref{fig:nivel2} presenta el desglose de los subprocesos internos de procesamiento de datos.

\begin{figure}[H]
    \centering
    \includegraphics[width=0.86\textwidth]{Figuras/nivel2.png} 
    \caption{Diagrama de Nivel 2: Mecánica interna del bloque de procesamiento (A2) (elaboración propia).}
    \label{fig:nivel2}
\end{figure}

\subsubsection{Subprocesos del Nivel 2}
Este nivel descompone el algoritmo de filtrado en cuatro tareas críticas:
\begin{itemize}
    \item \textbf{A21 (Lectura y amortiguación de entrada):} captura el flujo disponible en el puerto serial y en la interfaz de red, tolerando ráfagas y discontinuidades sin bloquear el proceso principal.
    \item \textbf{A22 (Clasificación temprana):} distingue si la sesión corresponde a línea base, telemetría de collar u ordeño completo, y separa el tratamiento de cada dominio antes de intentar cualquier correlación.
    \item \textbf{A23 (Normalización y metadatos):} transforma la evidencia aceptada en mensajes estructurados, con identificador de sesión, marca temporal, tipo de evento y resumen mínimo de red.
    \item \textbf{A24 (Persistencia y publicación):} decide si el lote puede publicarse, si debe quedar retenido por lista de permitidos o si debe almacenarse localmente para reenvío posterior.
\end{itemize}

\subsection{Lógica operativa de la pasarela y diagrama de flujo}
Para llevar la lógica anterior a la operación, se diseñó un flujo de control acotado y coherente con el estado actual del proyecto. En esta iteración, la pasarela no intenta decodificar todo el protocolo original del sistema en el borde; se comporta como un servicio determinista que clasifica el dominio de entrada, aplica reglas mínimas de seguridad, normaliza mensajes y decide si publica o retiene localmente la evidencia.

\colorlet{GatewayCoreBorder}{cyan!65!blue}
\colorlet{GatewayCoreFill}{cyan!10}
\colorlet{GatewayControlBorder}{teal!60!black}
\colorlet{GatewayControlFill}{teal!8}
\colorlet{GatewayExternalBorder}{gray!60}
\colorlet{GatewayExternalFill}{gray!8}
\colorlet{GatewayDecisionBorder}{orange!80!black}
\colorlet{GatewayDecisionFill}{orange!12}
\colorlet{GatewayAlertBorder}{red!70!black}
\colorlet{GatewayAlertFill}{red!8}
\colorlet{GatewayArrow}{black}
\colorlet{GatewayAux}{gray!60}

\begin{figure}[H]
    \centering
    \resizebox{!}{0.50\textheight}{%
    \begin{tikzpicture}[
        node distance = 1.3cm and 1.8cm,
        auto,
        font=\small,
        startEnd/.style = {rectangle, rounded corners, draw=GatewayControlBorder, thick, fill=GatewayControlFill, text width=4.1cm, align=center, minimum height=1cm},
        process/.style = {rectangle, rounded corners=3pt, draw=GatewayCoreBorder, thick, fill=GatewayCoreFill, text width=4.8cm, align=center, minimum height=1cm},
        io/.style = {rectangle, rounded corners=3pt, draw=GatewayExternalBorder, thick, fill=GatewayExternalFill, text width=4.8cm, align=center, minimum height=1cm},
        decision/.style = {diamond, draw=GatewayDecisionBorder, thick, fill=GatewayDecisionFill, text width=3.2cm, align=center, inner sep=0pt},
        db/.style = {cylinder, shape border rotate=90, draw=GatewayExternalBorder, thick, aspect=0.25, fill=GatewayExternalFill, text width=3.8cm, align=center, minimum height=1.5cm},
        arr/.style = {-Latex, thick, GatewayArrow}
    ]
    \node (A) [startEnd] {Inicio del servicio perimetral};
    \node (B) [process, below=of A] {Inicializar configuración, lista de permitidos y publicador};
    \node (C) [io, below=of B] {Leer sesión o lote de captura disponible};
    \node (D) [process, below=of C] {Clasificar dominio: línea base, collar u ordeño completo};
    \node (E) [decision, below=of D] {¿Política de red válida?};
    \node (F) [process, below=of E] {Normalizar a mensajes \ac{JSON} con tópico, tipo y metadatos};
    \node (G) [decision, below=of F] {¿Destino \ac{MQTT} disponible?};
    \node (H) [io, below=of G] {Publicar por \ac{MQTT} o espejo local};
    \node (I) [db, right=1.8cm of G] {Cola local \textit{JSONL}\\ almacenamiento y reenvío};
    \node (J) [process, right=1.8cm of E] {Retener lote y marcar observación de política};

    \draw [arr] (A) -- (B);
    \draw [arr] (B) -- (C);
    \draw [arr] (C) -- (D);
    \draw [arr] (D) -- (E);
    \draw [arr] (E) -- node[left] {Sí} (F);
    \draw [arr] (F) -- (G);
    \draw [arr] (G) -- node[left] {Sí} (H);
    \draw [arr] (G.east) -- node[above] {No} (I.west);
    \draw [arr] (E.east) -- node[above] {No} (J.west);
    \draw [arr] (H.west) -- ++(-1.0,0) |- (C.west);
    \draw [arr] (I.east) -- ++(1.5,0) |- (C.east);
    \draw [arr] (J.north) |- (C.east);
    \end{tikzpicture}
    }
    \caption{Flujo operativo de la primera iteración de la pasarela perimetral (elaboración propia).}
    \label{fig:flujo_python}
\end{figure}

\subsubsection{Criterios operativos y resiliencia de borde}
El flujo operativo de la Figura~\ref{fig:flujo_python} delimita qué debe resolver la pasarela y qué se reserva para el análisis posterior. En el borde se concentran la clasificación, la normalización y las decisiones mínimas de publicación; el análisis de mayor detalle se realiza después, con la evidencia ya consolidada. Bajo este criterio, el esquema se organiza en cuatro elementos principales:

\begin{table}[H]
\centering
\caption{Criterios operativos de resiliencia aplicados por la pasarela.}
\label{tab:criterios_resiliencia_borde}
\footnotesize
\setlength{\tabcolsep}{4pt}
\renewcommand{\arraystretch}{1.15}
\begin{tabularx}{\linewidth}{|>{\raggedright\arraybackslash}p{3.4cm}|>{\raggedright\arraybackslash}X|>{\raggedright\arraybackslash}X|}
\hline
\rowcolor{gray!20} \TableHeaderFirst{Elemento} & \TableHeader{Decisión en el borde} & \TableHeader{Efecto operativo} \\ \hline
Separación temprana por dominio & Distingue sesiones de línea base, telemetría de collar u ordeño completo antes de cualquier correlación. & Evita inferencias precipitadas y reduce la mezcla de significados antes del transporte. \\ \hline
Política mínima de seguridad & Valida IP y puerto esperados contra la lista de permitidos del dominio biótico. & Retiene los lotes fuera de política sin descartarlos y los deja señalados para revisión. \\ \hline
Normalización trazable & Encapsula tópico, tipo de evento, identificador de muestra y marca temporal en mensajes \acs{JSON}. & Permite trabajar con registros estructurados y reconstruir el origen de cada observación. \\ \hline
Persistencia y publicación diferida & Usa una cola \textit{JSONL} de almacenamiento y reenvío cuando el destino \ac{MQTT} no está disponible. & Desacopla la continuidad de captura de la disponibilidad inmediata de la red de salida. \\ \hline
\end{tabularx}
\end{table}

Con estos criterios, el modelo conceptual de amenazas del Capítulo~\ref{chapter:chapter03} se traduce en decisiones de diseño: proteger la disponibilidad e integridad de la telemetría, reducir el efecto de fallos no maliciosos en la \acs{LAN} y evitar publicaciones fuera de política. Por esa razón, los controles no se presentan como una capa aislada, sino como mecanismos operativos concretos: lista de permitidos, \acs{TLS}~1.3, autenticación mutua, registros locales y almacenamiento con reenvío.

\subsection{Implementación del servicio perimetral en Python}
La primera iteración de la pasarela perimetral se desarrolló en Python 3.13 para validar la ruta básica de normalización, política, persistencia y publicación. En esta etapa, el componente se ejecutó sobre sesiones ya procesadas para probar el contrato de datos y la lógica de borde antes de acoplarlo a captura en vivo. En lugar de incorporar el código fuente completo en el cuerpo del informe, la implementación se resume por su contrato operativo: directorio de sesión procesada como entrada, raíz de tópicos \ac{MQTT}, rutas locales de cola y espejo de publicados, modo de validación sin publicación externa y estado de publicación o retención de los mensajes.

Antes de conectar el nodo a un intermediario \ac{MQTT} en campo, la primera iteración se validó en aislamiento: el modo \texttt{dry\_run=True} desactiva la publicación externa y redirige la salida a archivos \texttt{.json} y \texttt{.jsonl} locales. Esta secuencia permitió comprobar que la cadena interna (normalización, política y contrato de salida) produjera salidas estables y revisables antes de depender de infraestructura de red activa.

En esta prueba se integraron configuración, política, normalización y salida en una misma ruta revisable. Con ello se confirmaron tres decisiones que se mantienen en las siguientes etapas: el contrato \ac{JSON} como formato de intercambio, la cola local \textit{JSONL} como mecanismo de persistencia y la publicación controlada, ya sea hacia el intermediario \ac{MQTT} o hacia un espejo local.

\subsubsection{Arquitectura modular del paquete \texttt{fincadiag.\allowbreak gateway}}
La Figura~\ref{fig:gateway_modules} presenta la organización interna del paquete \texttt{fincadiag.\allowbreak gateway}. Cada módulo asume una responsabilidad delimitada, lo que permite sustituir el publicador: puede usarse un espejo local de prueba o un cliente \ac{MQTT} real sin modificar el resto de la cadena.

\begin{figure}[H]
\centering
\resizebox{\linewidth}{!}{%
\begin{tikzpicture}[
    node distance=1.0cm and 1.6cm,
    font=\small,
    mod/.style={rectangle, rounded corners=3pt, draw=GatewayCoreBorder, thick, fill=GatewayCoreFill, text width=3.6cm, minimum height=1.1cm, align=center},
    cfg/.style={rectangle, rounded corners=3pt, draw=GatewayControlBorder, thick, fill=GatewayControlFill, text width=3.6cm, minimum height=1.1cm, align=center},
    ext/.style={rectangle, rounded corners=3pt, draw=GatewayExternalBorder, thick, fill=GatewayExternalFill, text width=3.0cm, minimum height=1.0cm, align=center, font=\small\itshape},
    arr/.style={-Latex, thick, GatewayArrow}
]

% Fila superior: config y runtime
\node[cfg] (config) {\textbf{config.py}\\ \scriptsize GatewayConfig\\ \scriptsize (TLS, lista de permitidos, rutas)};
\node[mod, right=2.2cm of config] (runtime) {\textbf{runtime.py}\\ \scriptsize GatewayRuntime\\ \scriptsize (orquestación)};

% Fila media: normalizer, policy
\node[mod, below left=1.4cm and 0.2cm of runtime] (normalizer) {\textbf{normalizer.py}\\ \scriptsize JSON por dominio\\ \scriptsize (sesión, línea base, pcap, collar, correlación)};
\node[mod, below right=1.4cm and 0.2cm of runtime] (policy) {\textbf{policy.py}\\ \scriptsize AllowlistPolicy\\ \scriptsize (IP + puerto 6001)};

% Fila inferior: publisher, store, models
\node[mod, below=1.4cm of normalizer] (publisher) {\textbf{publisher.py}\\ \scriptsize PahoPublisher /\\ \scriptsize FileMirrorPublisher};
\node[mod, below=1.4cm of policy] (store) {\textbf{store.py}\\ \scriptsize JsonlSpoolStore\\ \scriptsize (almacenamiento y reenvío)};
\node[cfg, below right=0.6cm and 0.0cm of publisher] (models) {\textbf{models.py}\\ \scriptsize GatewayMessage\\ \scriptsize PublicationResult};

% Externos
\node[ext, left=1.4cm of publisher] (broker) {Intermediario MQTT\\ \scriptsize (Mosquitto + TLS\,1.3)};
\node[ext, right=1.4cm of store] (spool) {Cola JSONL\\ \scriptsize (disco local)};

% Flechas
\coordinate (runtime_hub) at ($(runtime.south)+(0,-1.35cm)$);
\coordinate (publisher_runtime_in) at ([yshift=0.18cm]publisher.east);
\coordinate (publisher_normalizer_in) at (publisher.north);
\coordinate (publisher_models_in) at ([xshift=0.18cm]publisher.south);
\coordinate (store_runtime_in) at ([yshift=0.18cm]store.west);
\coordinate (store_policy_in) at (store.north);
\coordinate (store_models_in) at ([xshift=-0.18cm]store.south);
\coordinate (runtime_pub_lane) at ([xshift=0.65cm]publisher_runtime_in);
\coordinate (runtime_store_lane) at ([xshift=-0.65cm]store_runtime_in);
\coordinate (runtime_store_out) at ([xshift=0.22cm]runtime.south);
\coordinate (runtime_store_hub) at ($(runtime_store_out)+(0,-1.35cm)$);
\draw[arr] (config) -- (runtime);
\draw[arr] (runtime) -- (normalizer);
\draw[arr] (runtime) -- (policy);
\draw[arr] (runtime.south) -- (runtime_hub) -| (runtime_pub_lane) -- (publisher_runtime_in);
\draw[{Latex}-{Latex}, thick, GatewayArrow] (runtime_store_out) -- (runtime_store_hub) -| (runtime_store_lane) -- (store_runtime_in);
\draw[arr] (normalizer.south) -- (publisher_normalizer_in);
\draw[arr] (policy.south) -- (store_policy_in);
\draw[arr] (publisher) -- (broker);
\draw[arr] (store) -- (spool);
\draw[arr, dashed, GatewayAux] (models.west) -| (publisher_models_in);
\draw[arr, dashed, GatewayAux] (models.east) -| (store_models_in);

\end{tikzpicture}
}
\caption{Arquitectura modular del paquete \texttt{fincadiag.\allowbreak gateway}: relación entre configuración, normalización, política, publicación y persistencia local (elaboración propia).}
\label{fig:gateway_modules}
\end{figure}

En esta arquitectura, \texttt{runtime.py} actúa como orquestador central: recibe una sesión procesada, la normaliza en mensajes \ac{JSON}, evalúa la política de red y decide si publica al intermediario \ac{MQTT} o retiene el lote en la cola local. Los módulos \texttt{models.py} y \texttt{config.py} proporcionan las estructuras de datos compartidas, mientras que la capa de publicación admite dos implementaciones intercambiables: \texttt{PahoPublisher} para entornos con sincronización \ac{MQTT} y soporte de \ac{TLS}~1.3, y \texttt{FileMirrorPublisher} para validación sin publicación externa.

La relación entre \texttt{runtime.py} y \texttt{store.py} es bidireccional. En condiciones normales, \texttt{runtime.py} escribe en la cola local cuando la publicación falla; en el modo \texttt{--watch}, el método \texttt{drain\_spool()} lee los lotes pendientes desde el disco y reintenta su entrega contra el intermediario \ac{MQTT}. Una vez confirmada la entrega, el archivo \textit{JSONL} se elimina de la cola, cerrando el ciclo de almacenamiento y reenvío sin intervención manual.


Una vez revisado que la cadena interna producía mensajes coherentes, el siguiente paso fue probarla con sesiones reales ya procesadas. El objetivo fue comprobar que funcionara de manera estable, conservara la cobertura por dominio y mantuviera coherencia semántica antes de conectarla a captura en vivo.

\subsubsection{Verificación fuera de línea y ajustes de la primera iteración}
La verificación fuera de línea de la pasarela, en ejecución sin publicación externa, se utilizó para revisar estabilidad operativa, coherencia semántica y cobertura de mensajes sobre sesiones procesadas con distinto nivel de detalle analítico. En esta fase también se detectó y corrigió un error en el resumen de correlación, y se añadieron apoyos de revisión como el ajuste de \texttt{run\_gateway.bat}, la salida adicional \texttt{.readable.json} y el reporte explícito de rutas de salida. La Tabla~\ref{tab:gateway_offline_validation_cases} sintetiza esta revisión como control de implementación: indica qué aspecto se revisó, qué evidencia se utilizó y qué decisión técnica quedó incorporada, sin trasladar a este capítulo los resultados cuantitativos que se discuten en el Capítulo~\ref{chapter:chapter06}.

\begin{table}[H]
\centering
\caption{Síntesis metodológica de la verificación fuera de línea de la pasarela.}
\label{tab:gateway_offline_validation_cases}
\footnotesize
\setlength{\tabcolsep}{3.5pt}
\renewcommand{\arraystretch}{1.05}
\begin{tabularx}{0.98\linewidth}{|>{\raggedright\arraybackslash}p{2.7cm}|>{\raggedright\arraybackslash}X|>{\raggedright\arraybackslash}X|}
\hline
\rowcolor{gray!20} \TableHeaderFirst{Aspecto revisado} & \TableHeader{Evidencia utilizada} & \TableHeader{Decisión o ajuste derivado} \\ \hline
Estructura de mensajes & Sesiones procesadas con diferente volumen de eventos, alertas y registros por vaca. & Se confirmó el contrato \ac{JSON} y la separación por tipo de mensaje antes de publicar. \\ \hline
Política de publicación & Ejecuciones sin publicación externa, con espejo local y cola \textit{JSONL}. & Se mantuvo el modo de revisión local y el reporte explícito de rutas de salida. \\ \hline
Coherencia semántica & Casos con soporte de correlación disponible y distintos niveles de observación de campo. & Se corrigió el conteo de correlación y se evitó forzar causalidad sin soporte simultáneo. \\ \hline
Preparación para despliegue & Lotes representativos de sesiones reales ya consolidadas por el motor de análisis. & Se seleccionó la cadena normalización--política--persistencia--publicación para la preintegración con \ac{MQTT}/\ac{TLS}. \\ \hline
\end{tabularx}
\end{table}

\subsubsection{Preintegración controlada con intermediario \acs{MQTT} y \acs{TLS} de laboratorio}

Para verificar que la pasarela podía publicar sobre un canal cifrado real, y no solo contra un espejo local, se montó un intermediario \ac{MQTT} Mosquitto en una estación de laboratorio, con \ac{TLS}~1.3 y autenticación mutua mediante certificados. Se generaron una autoridad certificadora local, un certificado de servidor con \texttt{SAN=localhost} y un certificado de cliente, todos firmados por la misma autoridad. La prueba se ejecutó sobre una sesión combinada ya procesada, apuntando a \texttt{localhost:8883}.

El ensayo confirmó que la cadena completa —normalización, política, publicación \acs{MQTT} con \acs{TLS}~1.3 y espejo local— podía operar de extremo a extremo sin intervención manual. Los conteos de publicación y retención se reservan para el Capítulo~\ref{chapter:chapter06}, donde se interpretan como parte de la validación.

También se ejecutó una prueba de resiliencia con el intermediario \ac{MQTT} intencionalmente detenido. En la versión inicial, la pasarela terminaba con una excepción de conexión (\texttt{ConnectionRefusedError}) y no conservaba el lote pendiente. Para corregir ese problema, se añadió un bloque \texttt{try/except} en \texttt{runtime.py} que, ante fallo de publicación, preserva el lote completo en la cola local \textit{JSONL}. Tras el ajuste, la prueba confirmó que los mensajes quedaban retenidos sin pérdida de datos ni terminación abrupta; los conteos detallados se trasladan al capítulo de resultados.

Con estas pruebas se cerró la verificación de laboratorio local. Los resultados del despliegue en campo sobre la Raspberry Pi (\texttt{gateway-esmeralda}), incluyendo el modo residente continuo \texttt{--watch}, se presentan en el Capítulo~\ref{chapter:chapter06}.

Con la arquitectura del paquete ya definida y probada en laboratorio, el siguiente paso consistió en instalarla sobre un nodo de campo estable, de bajo consumo y capaz de operar de forma desatendida. La Raspberry Pi se seleccionó como plataforma de despliegue por su balance entre recursos de procesamiento, soporte nativo de interfaces seriales y red, y la posibilidad de ejecutar un sistema operativo liviano sin interfaz gráfica.

\subsection{Configuración del nodo perimetral (Raspberry Pi OS Lite y systemd)}
Se seleccionó \textbf{Raspberry Pi OS Lite (64-bit)} para favorecer la estabilidad del nodo y reducir procesos no esenciales. Al tratarse de una variante sin interfaz gráfica, esta elección libera recursos para la adquisición serial y el transporte. La persistencia del servicio se formalizó mediante una unidad de \textbf{systemd} con reinicio automático y retardo controlado de reintento.

La transición operativa del despliegue se resume para evitar ambigüedades con la metodología del Capítulo~\ref{chapter:chapter04}. Hasta el \textbf{19 de marzo}, las capturas se ejecutaron principalmente como corridas manuales controladas desde una computadora con Windows; esas corridas permitieron explorar el canal, ajustar el entorno y depurar el instrumento. A partir del \textbf{20 de marzo}, la operación migró a Linux y se calendarizó mediante \texttt{cron}. El servicio residente con \textit{systemd} corresponde a la etapa de intervención consolidada descrita para el 11 de mayo; por tanto, no se confunde con la calendarización previa, sino que la reemplaza como mecanismo de continuidad del nodo perimetral.


\begin{table}[H]
\centering
\caption{Resumen de configuración operativa del servicio residente.}
\label{tab:systemd_gateway_summary}
\footnotesize
\setlength{\tabcolsep}{3.5pt}
\renewcommand{\arraystretch}{1.0}
\begin{tabularx}{\linewidth}{|>{\raggedright\arraybackslash}p{2.1cm}|>{\raggedright\arraybackslash}X|>{\raggedright\arraybackslash}p{4.4cm}|}
\hline
\rowcolor{gray!20} \TableHeaderFirst{Bloque} & \TableHeader{Configuración aplicada} & \TableHeader{Función en la pasarela} \\ \hline
Arranque & Espera de red antes de iniciar el servicio. & Evitar fallos por arranque prematuro. \\ \hline
Proceso & Módulo residente \texttt{fincadiag.gateway.runtime}. & Mantener la pasarela como servicio desatendido. \\ \hline
Directorios & Separación entre sesiones procesadas, cola local y espejo de publicados. & Prevenir reprocesamiento y habilitar cola/reenvío. \\ \hline
Transporte & Broker local, puerto seguro y certificados de cliente. & Publicar por \ac{MQTT}/\ac{TLS} dentro del perímetro. \\ \hline
Continuidad & Modo \texttt{-{}-watch}, intervalo de 60\,s y reinicio automático. & Drenar cola, detectar sesiones y recuperarse ante fallos. \\ \hline
\end{tabularx}
\end{table}

La Tabla~\ref{tab:systemd_gateway_summary} resume la configuración del servicio, pero su lectura operativa requiere explicitar el ciclo que ejecuta el nodo una vez iniciado. En este punto, la pasarela deja de entenderse solo como un conjunto de parámetros de \textit{systemd} y pasa a comportarse como un proceso residente con una secuencia repetitiva de recuperación, búsqueda y publicación. La política \texttt{Restart=always}, junto con el retardo \texttt{RestartSec=5}, permite que \textit{systemd} relance el servicio ante una interrupción, mientras que el modo \texttt{--watch} mantiene el proceso activo. En cada ciclo de 60 segundos, el servicio drena primero la cola local de mensajes pendientes y luego revisa \texttt{--visit-dir} en busca de sesiones nuevas; una sesión se considera pendiente cuando aún no existe su archivo espejo en \texttt{--published-dir}. La Figura~\ref{fig:watch_cycle} representa esta lógica de operación continua.

\begin{figure}[H]
\centering
\resizebox{0.4\linewidth}{!}{%
\begin{tikzpicture}[
    node distance=0.85cm,
    font=\small,
    ctrl/.style={rectangle, rounded corners=6pt, draw=ThesisGreen, very thick, fill=ThesisGreen!8, text width=4.8cm, minimum height=1.0cm, align=center},
    mod/.style={rectangle, rounded corners=3pt, draw=GatewayCoreBorder, thick, fill=GatewayCoreFill, text width=4.8cm, minimum height=1.05cm, align=center},
    ext/.style={rectangle, rounded corners=3pt, draw=GatewayExternalBorder, thick, fill=GatewayExternalFill, text width=4.8cm, minimum height=1.05cm, align=center},
    arr/.style={-Latex, thick, GatewayArrow}
]
\node[ctrl] (A) {Inicio del ciclo \texttt{-{}-watch}};
\node[mod, below=of A] (B) {\texttt{drain\_spool()}\\\scriptsize Leer \texttt{spool/*.jsonl},\ publicar\\\scriptsize y eliminar al confirmar entrega};
\node[ext, below=of B] (C) {Escanear \texttt{-{}-visit-dir}\\\scriptsize Obtener sesiones ausentes\\\scriptsize en \texttt{published\_dir}};
\node[mod, below=of C] (D) {\texttt{publish\_session\_dir()}\\\scriptsize Para cada sesión nueva detectada};
\node[ctrl, below=of D] (E) {\texttt{time.sleep(60\,s)}\\\scriptsize Esperar antes del próximo ciclo};
\draw[arr] (A) -- (B);
\draw[arr] (B) -- (C);
\draw[arr] (C) -- (D);
\draw[arr] (D) -- (E);
\draw[arr] (E.west) -- ++(-2.0,0) |- (A.west);
\end{tikzpicture}
}
\caption{Ciclo de operación continua del modo \texttt{-{}-watch}: drenado de la cola local y detección de sesiones nuevas en cada iteración (elaboración propia).}
\label{fig:watch_cycle}
\end{figure}

Con el nodo configurado y el servicio residente, el paso siguiente es definir cómo se orquestan las sesiones de captura y qué instrumento realiza la adquisición en campo.

\subsection[Orquestación de captura]{Orquestación de captura: el módulo\\\texttt{FincaScheduler.py}}
Hasta este punto, el capítulo ha centrado la atención en la pasarela perimetral como componente de captura y publicación. Sin embargo, la pasarela no opera aislada: forma parte de un instrumento más amplio, denominado \textbf{FincaDiag}, que comprende tanto la adquisición de evidencia en campo como su posterior ordenamiento analítico. Dentro de este instrumento, el módulo \texttt{FincaScheduler} (v3.0) actúa como orquestador de las sesiones de captura, gestionando cada bloque como una secuencia de fases con estado persistente. La organización temporal de estos bloques, incluyendo horarios, fases y modos de captura, fue establecida metodológicamente en el Capítulo~\ref{chapter:chapter04} (Tabla~\ref{tab:calendarizacion_bloques}).

\subsubsection{Máquina de estados del orquestador}
El orquestador gestiona el ciclo de vida de cada sesión de captura mediante una máquina de estados que conserva su contexto en disco. La Figura~\ref{fig:scheduler_state_machine} formaliza las transiciones posibles, incluyendo las rutas de recuperación ante reinicio del sistema operativo o caída del proceso.

\begin{figure}[H]
\centering
\resizebox{0.93\linewidth}{!}{%
\begin{tikzpicture}[
    node distance=1.6cm and 2.7cm,
    font=\small,
    state/.style={circle, draw=GatewayCoreBorder, thick, fill=GatewayCoreFill, text width=2.35cm, minimum size=2.75cm, inner sep=1.5pt, align=center},
    stateok/.style={circle, draw=GatewayControlBorder, thick, fill=GatewayControlFill, text width=2.35cm, minimum size=2.75cm, inner sep=1.5pt, align=center},
    statefail/.style={circle, draw=GatewayAlertBorder, thick, fill=GatewayAlertFill, text width=2.35cm, minimum size=2.75cm, inner sep=1.5pt, align=center},
    statewarn/.style={circle, draw=GatewayDecisionBorder, thick, fill=GatewayDecisionFill, text width=2.35cm, minimum size=2.75cm, inner sep=1.5pt, align=center},
    init/.style={circle, fill=GatewayArrow, minimum size=6pt, inner sep=0pt},
    trans/.style={-Latex, thick, GatewayArrow},
    recovery/.style={-Latex, thick, dashed, GatewayAlertBorder},
    edgelabel/.style={font=\scriptsize, fill=white, rounded corners=1pt, inner sep=1.8pt, outer sep=1.0pt, align=center}
]

% Estados
\node[init] (init) {};
\node[state, right=1.8cm of init] (idle) {Descanso\\ \scriptsize (entre bloques)};
\node[state, right=of idle] (pending) {Pendiente\\ \scriptsize (bloque detectado)};
\node[state, right=of pending] (running) {Ejecutando\\ \scriptsize (fase activa)};
\node[stateok, below right=1.65cm and 0.0cm of running] (completed) {Completado\\ \scriptsize (fases correctas)};
\node[statewarn, below left=1.65cm and 0.0cm of running] (incomplete) {Incompleto\\ \scriptsize (tiempo agotado)};
\node[statefail, above=1.6cm of running] (crash) {Fallo detectado\\ \scriptsize (reinicio / excepción)};

\coordinate (completedreturn) at ($(completed.south)+(0,-0.72)$);
\coordinate (idlereturnlong) at (idle.south |- completedreturn);
\coordinate (incompletereturn) at ($(incomplete.south)+(0,-0.42)$);
\coordinate (idlereturnshort) at (idle.south |- incompletereturn);

% Transiciones normales
\draw[trans] (init) -- node[edgelabel, pos=0.52, above=8pt] {arranque} node[edgelabel, pos=0.52, below=8pt] {cargar agenda} (idle);
\draw[trans] (idle) -- node[edgelabel, midway, above=8pt] {hora de inicio} node[edgelabel, midway, below=8pt] {[bloque programado]} (pending);
\draw[trans] (pending) -- node[edgelabel, midway, above=8pt] {bloque válido} node[edgelabel, midway, below=8pt] {crear estado} (running);
\draw[trans] (running.east) .. controls ++(1.45,-0.78) and ++(1.45,0.78) .. node[edgelabel, pos=0.50, right=6pt] {fase correcta [quedan más] / avanzar} (running.east);
\draw[trans] (running.south east) to[out=-35,in=105] (completed.north);
\node[edgelabel] at ($(running.south east)!0.50!(completed.north) + (1.04,0.46)$) {fase correcta [última]};
\draw[trans] (running.south west) to[out=-145,in=75] (incomplete.north);
\node[edgelabel] at ($(running.south west)!0.50!(incomplete.north) + (-0.48,0.46)$) {tiempo agotado};
\draw[trans] (completed.south) -- (completedreturn) -- (idlereturnlong) -- (idle.south);
\node[edgelabel, rotate=90] at ($(idlereturnlong)!0.45!(idle.south) + (-0.22,-0.18)$) {nuevo bloque / limpiar};
\draw[trans] (incomplete.south) -- (incompletereturn) -- (idlereturnshort) -- (idle.south);

% Transiciones de recuperación
\draw[recovery] (running.north west) to[out=125,in=-125] node[edgelabel, pos=0.50, left=2pt] {fallo detectado / registrar} (crash.south west);
\draw[recovery] (crash.south east) to[out=-55,in=55] node[edgelabel, pos=0.50, right=2pt] {reset\_if\_needed / reanudar} (running.north east);

% Heartbeat
\node[font=\tiny\itshape, text=black!55, right=0.3cm of crash] {pulso cada 15\,s};

\end{tikzpicture}
}
\caption{Máquina de estados del orquestador \texttt{FincaScheduler}: transiciones nominales y rutas de recuperación ante fallos (elaboración propia).}
\label{fig:scheduler_state_machine}
\end{figure}

Para resistir cortes de energía, el estado se guarda en disco mediante una escritura atómica (\texttt{os.replace}), de modo que el archivo siempre quede en un estado válido: o se conserva íntegra la versión anterior o se sustituye por completo por la nueva, sin dejar registros parciales. Al reiniciarse, el orquestador compara el \texttt{boot\_id} del sistema con el valor registrado en el estado: si ambos difieren, interpreta que hubo un reinicio y reanuda la fase interrumpida desde el punto en que se detuvo, descontando los segundos ya ejecutados. Esta estrategia contribuye a mantener el \ac{MTTR} dentro del umbral de diseño ($\leq 60\,\text{s}$) y a minimizar los vacíos de captura durante incidencias del suministro eléctrico, frecuentes en el entorno rural de la finca.

El orquestador determina cuándo iniciar una sesión de captura y la mantiene bajo control durante su ejecución, pero la adquisición propiamente dicha (tráfico serial, paquetes de red y línea base) corre a cargo del programa de campo que se ejecuta directamente sobre el nodo perimetral.

\subsection{Instrumento de adquisici\'{o}n de campo: \texttt{FincaDiag.py}}
\texttt{FincaDiag.py} es el ejecutable principal de adquisici\'{o}n desplegado en la Raspberry Pi. Sus primeras pruebas operativas iniciaron el \textbf{23 de febrero} en el laboratorio local; posteriormente, desde el \textbf{20 de marzo}, se consolid\'{o} como instrumento principal de captura en campo. La versi\'{o}n inicial evolucion\'{o} conforme se obtuvo mayor informaci\'{o}n sobre el funcionamiento del ecosistema en estudio. Su funci\'{o}n no es procesar visitas ya capturadas, sino realizar la captura en vivo de evidencia: l\'{i}nea base de red, tr\'{a}fico serial, paquetes \acs{UDP} del collar o una combinaci\'{o}n de ellos, seg\'{u}n el modo seleccionado. La Tabla~\ref{tab:fincadiag_modos} describe la funci\'{o}n asociada a cada valor del argumento \texttt{-m}.

\begin{table}[H]
\centering
\caption{Modos de adquisici\'{o}n disponibles en \texttt{FincaDiag.py}.}
\label{tab:fincadiag_modos}
\footnotesize
\renewcommand{\arraystretch}{1.10}
\begin{tabular}{|c|l|}
\hline
\rowcolor{gray!20} \TableHeaderFirst{Modo} & \TableHeader{Descripci\'{o}n} \\ \hline
\texttt{-m 1} & Antena \acs{UDP} + \acs{PCAP} filtrado en puerto 6001 \\ \hline
\texttt{-m 2} & Serial + \acs{PCAP} completo (bloques de orde\~{n}o) \\ \hline
\texttt{-m 3} & Solo \acs{PCAP} completo (bloques normales) \\ \hline
\texttt{-m 4} & Línea base de red \\ \hline
\texttt{-m 5} & Serial + antena \acs{UDP} + \acs{PCAP} en paralelo \\ \hline
\end{tabular}
\end{table}

Con esta clasificaci\'{o}n, la operación de campo se resume en tres argumentos: el modo de adquisición, la duración de la ventana y los permisos de ejecución. En el modo combinado \texttt{-m 5}, \texttt{FincaDiag.py} captura simultáneamente serial, antena \acs{UDP} y \acs{PCAP}, de acuerdo con la Tabla~\ref{tab:fincadiag_modos}; una ventana típica de \texttt{-t 300} fija 300 segundos de captura. La ejecución con privilegios elevados es necesaria porque el programa accede a \texttt{tcpdump}, a la interfaz de red y al puerto serial.

Con estos par\'{a}metros definidos, la Figura~\ref{fig:fincadiag_pipeline} describe la secuencia operativa real del programa: validaci\'{o}n de permisos, selecci\'{o}n del modo de adquisici\'{o}n indicado en la Tabla~\ref{tab:fincadiag_modos}, configuraci\'{o}n de la duraci\'{o}n, preparaci\'{o}n de artefactos de salida, apertura de los canales de captura, bucle de lectura durante la ventana de captura y escritura del registro de cierre al finalizar.


\begin{figure}[H]
\centering
\resizebox{0.88\linewidth}{!}{%
\begin{tikzpicture}[
    node distance=0.72cm and 1.7cm,
    transform shape,
    font=\small,
    phase/.style={rectangle, rounded corners=4pt, draw=GatewayCoreBorder, thick, fill=GatewayCoreFill, text width=4.15cm, minimum height=0.88cm, align=center},
    io/.style={rectangle, rounded corners=4pt, draw=GatewayExternalBorder, thick, fill=GatewayExternalFill, text width=4.15cm, minimum height=0.88cm, align=center},
    decision/.style={diamond, aspect=2.1, draw=GatewayDecisionBorder, thick, fill=GatewayDecisionFill, text width=2.75cm, align=center, inner sep=1pt, font=\small},
    junction/.style={circle, fill=GatewayArrow, inner sep=1.4pt},
    lane/.style={font=\scriptsize\bfseries, text=black!65},
    arrow/.style={-Latex, thick, GatewayArrow}
]

\node[io] (start) {Inicio\\\texttt{sudo python3 FincaDiag.py}};
\node[phase, below=of start] (root) {Validar permisos root y leer argumentos};
\node[decision, below=0.75cm of root] (mode) {¿Modo N°?};

\node[lane, below left=0.45cm and 2.4cm of mode] (baselineTitle) {};
\node[phase, below=0.6cm of baselineTitle] (baseline) {Generar línea base de red};
\node[io, below=of baseline] (baselineOut) {Registros de línea base};

\node[lane, below right=0.45cm and 2.4cm of mode] (captureTitle) {};
\node[decision, below=0.6cm of captureTitle] (time) {¿Tiempo definido?};
\node[phase, below=of time] (prep) {Crear directorio y registros de salida};
\node[phase, below=of prep] (tcp) {Iniciar \texttt{tcpdump}\\PCAP completo y filtrado};
\node[decision, below left=0.50cm and 1.15cm of tcp] (udp) {¿UDP activo?};
\node[decision, below right=0.50cm and 1.15cm of tcp] (serial) {¿Serial activo?};
\node[junction, below=0.95cm of tcp] (captureJoin) {};
\coordinate (captureDecisionLabelLevel) at ($(captureJoin)+(0,0.30cm)$);
\node[font=\scriptsize, inner sep=0pt] at (captureDecisionLabelLevel) {S\'{i}};
\node[io, text width=2.10cm, minimum height=0.70cm, font=\scriptsize, below left=0.62cm and 0.28cm of udp] (udpArtifact) {Registros de captura};
\node[io, text width=2.10cm, minimum height=0.70cm, font=\scriptsize, below right=0.62cm and 0.28cm of serial] (serialArtifact) {Registros de captura};
\node[phase, below=1.18cm of captureJoin] (loop) {Bucle de captura durante la ventana definida};
\node[phase, below=of loop] (finish) {Finalizar captura y escribir\\registro JSON de cierre};
\node[io, below=of finish] (captureOut) {Registros de captura en disco};

\draw[arrow] (start) -- (root);
\draw[arrow] (root) -- (mode);
\draw[arrow] (mode.west) -| node[pos=0.25, above=0.10cm, font=\scriptsize, inner sep=0pt] {Modo 4} (baseline.north);
\draw[arrow] (baseline) -- (baselineOut);
\draw[arrow] (mode.east) -| node[pos=0.25, above=0.10cm, font=\scriptsize, inner sep=0pt] {Modos 1--3, 5} (time.north);
\draw[arrow] (prep.west) -- ++(-0.60,0) |- node[pos=0.72, above=0.16cm, font=\scriptsize, inner sep=0pt] {Línea base inicial} (baseline.east);
\draw[arrow] (time) -- (prep);
\node[font=\scriptsize, inner sep=0pt] at ($($(time)!0.5!(prep)$)+(0.18cm,0)$) {S\'{i}};
\coordinate (timeNoGuide) at ($(time.east)+(1.05,0)$);
\draw[arrow] (time.east) -- (timeNoGuide) |- (start.east);
\node[font=\scriptsize, inner sep=0pt] at ($($(time.east)!0.5!(timeNoGuide)$)+(0,0.24cm)$) {No};
\draw[arrow] (prep) -- (tcp);
\draw[arrow] (tcp) -| (udp);
\draw[arrow] (tcp) -| (serial);
\draw[arrow] (udp.east) -- ++(0.35,0) |- (captureJoin.west);
\draw[arrow] (serial.west) -- ++(-0.35,0) |- (captureJoin.east);
\draw[arrow] (udp.west) -- ++(-0.34,0) -| (udpArtifact.north);
\draw[arrow] (serial.east) -- ++(0.34,0) -| (serialArtifact.north);
\node[font=\scriptsize, inner sep=0pt] at ($($(udp.west)!0.58!($(udp.west)+(-0.34,0)$)$)+(0,0.14cm)$) {No};
\node[font=\scriptsize, inner sep=0pt] at ($($(serial.east)!0.58!($(serial.east)+(0.34,0)$)$)+(0,0.24cm)$) {No};
\draw[arrow] (captureJoin) -- (loop);
\draw[arrow] (loop) -- (finish);
\draw[arrow] (finish) -- (captureOut);
\coordinate (baselinePostGuide) at ($(baseline.west)+(-0.85,0)$);
\coordinate (baselinePostTurn) at (finish.west -| baselinePostGuide);
\coordinate (baselinePostEnd) at (baseline.west);
\draw[arrow] (finish.west) -- (baselinePostTurn) -- node[midway, left=0.16cm, font=\scriptsize, inner sep=0pt] {Línea base final} (baselinePostGuide) -- (baselinePostEnd);

\end{tikzpicture}%
}
\caption{Secuencia operativa de \texttt{FincaDiag.py}: validaci\'{o}n, selecci\'{o}n de modo, captura en vivo y escritura de registros en la Raspberry Pi (elaboraci\'{o}n propia).}
\label{fig:fincadiag_pipeline}
\end{figure}

De esta forma, \texttt{FincaDiag.py} realiza la adquisici\'{o}n de evidencia en campo y genera los insumos brutos que posteriormente alimentan tanto el motor de an\'{a}lisis como la pasarela perimetral. A partir de esa captura, la siguiente tarea de la pasarela consiste en convertir la evidencia admitida por la pol\'{i}tica de borde en mensajes normalizados y transportables. A continuación se describe el formato de salida y los mecanismos que lo preservan ante fallos de conectividad o reinicios del nodo.

\subsection{Salida de la pasarela: estructuración del mensaje \acs{JSON}}
La salida del servicio perimetral funciona como punto de transición entre la captura bruta y los módulos que reciben esos datos. En lugar de publicar archivos completos o un único mensaje biótico agregado, la pasarela transforma la evidencia admitida por la política de borde en mensajes \acs{JSON} ligeros, trazables y separados por tópico. Este formato facilita la interoperabilidad con el intermediario \ac{MQTT} Mosquitto, conserva el dominio de origen de cada observación y evita que las etapas posteriores dependan del detalle crudo de la captura.

En la versión implementada, cada mensaje conserva el contexto de sesión, el tópico de publicación y el tipo de evento que representa. De esta forma, quien lo reciba puede distinguir un resumen de sesión, una alerta, una correlación de ordeño o un evento individual sin reabrir los archivos fuente. La Tabla~\ref{tab:gateway_json_schema} resume los campos que componen el contrato de salida sin introducir un listado extenso en el cuerpo del capítulo.

\begin{table}[!htbp]
\centering
\caption{Campos principales del mensaje normalizado por la pasarela.}
\label{tab:gateway_json_schema}
\footnotesize
\renewcommand{\arraystretch}{1.12}
\begin{tabularx}{\linewidth}{|>{\raggedright\arraybackslash}p{3.2cm}|>{\raggedright\arraybackslash}X|>{\raggedright\arraybackslash}p{3.0cm}|}
\hline
\rowcolor{gray!20} \TableHeaderFirst{Campo} & \TableHeader{Contenido} & \TableHeader{Función} \\ \hline
\texttt{topic} & Ruta semántica con finca, visita, captura y tipo de evento. & Enrutamiento \ac{MQTT}. \\ \hline
\texttt{payload} & Métricas o datos normalizados del evento, por ejemplo eficiencia, desfase o conteos. & Carga útil analítica. \\ \hline
\texttt{qos} y \texttt{retain} & Parámetros de entrega y retención del intermediario. & Control de publicación. \\ \hline
\texttt{source\_sample\_id} & Identificador de la captura que originó el mensaje. & Trazabilidad de origen. \\ \hline
\texttt{event\_type} y \texttt{event\_timestamp} & Tipo lógico del evento y marca temporal normalizada. & Interpretación y ordenamiento. \\ \hline
\end{tabularx}
\end{table}

La estructura ilustrada favorece un crecimiento ordenado del sistema: el borde publica mensajes homogéneos, el intermediario \ac{MQTT} conserva la separación semántica por tópico y los módulos posteriores pueden añadir reglas de negocio sin reescribir la lógica de captura. Para el proyecto, esta decisión también aporta una ventaja metodológica, porque permite evaluar la estructura de datos de la pasarela antes de completar el despliegue en campo.

Una vez definido qué se publica y cómo se distingue cada evento, el siguiente problema operativo es garantizar que esos mensajes no se pierdan cuando la cadena de publicación no esté disponible. Por esta razón, la continuidad del servicio se aborda como una extensión directa del formato de salida: si el intermediario \ac{MQTT} falla o el nodo se reinicia, la pasarela debe retener los lotes ya normalizados y reanudarlos de forma controlada.

\subsection{Continuidad operativa de la pasarela: recuperación automática y tolerancia a fallos}
Para llevar esta política a la operación del nodo, la pasarela se apoya en dos mecanismos complementarios: el reinicio automático del servicio y la persistencia local de los lotes pendientes. El primero permite que el proceso vuelva a levantarse después de una excepción o de un reinicio del sistema; el segundo evita descartar mensajes ya normalizados cuando el intermediario \ac{MQTT} o la red no estén disponibles. La Figura~\ref{fig:flujo_rollback} organiza estas decisiones como un flujo de publicación, retención y recuperación.

\begin{figure}[H]
    \centering
    \makebox[\linewidth][c]{\hspace*{-1.2cm}\resizebox{!}{0.52\textheight}{%
    \begin{tikzpicture}[
        node distance = 1.4cm and 1.8cm,
        auto,
        font=\small,
        startEnd/.style = {rectangle, rounded corners, draw=GatewayControlBorder, thick, fill=GatewayControlFill, text width=4.0cm, align=center, minimum height=1cm},
        process/.style = {rectangle, rounded corners=3pt, draw=GatewayCoreBorder, thick, fill=GatewayCoreFill, text width=4.8cm, align=center, minimum height=1cm},
        io/.style = {rectangle, rounded corners=3pt, draw=GatewayExternalBorder, thick, fill=GatewayExternalFill, text width=4.8cm, align=center, minimum height=1cm},
        decision/.style = {diamond, draw=GatewayDecisionBorder, thick, fill=GatewayDecisionFill, text width=3.2cm, align=center, inner sep=0pt},
        alert/.style = {rectangle, rounded corners=3pt, draw=GatewayAlertBorder, thick, fill=GatewayAlertFill, text width=4.5cm, align=center, minimum height=1cm},
        arr/.style = {-Latex, thick, GatewayArrow}
    ]
    \node (A) [startEnd] {Inicio / reinicio por systemd};
    \node (B) [process, below=of A] {Cargar configuración y reabrir cola local};
    \node (C) [io, below=of B] {Procesar lote o sesión disponible};
    \node (D) [decision, below=of C] {¿Publicación correcta?};
    \node (E) [process, below left=0.85cm and 2.4cm of D] {Confirmar envío y espejar lote publicado};
    \node (F) [alert, below right=0.85cm and 2.4cm of D] {Retener en cola JSONL y marcar reintento};
    \node (G) [decision, below=2.8cm of D] {¿Servicio caído o excepción?};
    \node (H) [process, below left=1.4cm and 2.4cm of G] {Continuar flujo nominal};
    \node (I) [alert, below right=1.4cm and 2.4cm of G] {Reinicio automático y recuperación de la cola};
    \coordinate (serviceJoin) at ($(G.north)+(0,0.55cm)$);

    \draw [arr] (A) -- (B);
    \draw [arr] (B) -- (C);
    \draw [arr] (C) -- (D);
    \draw [arr] (D.west) -| node[pos=0.25, above] {Sí} (E.north);
    \draw [arr] (D.east) -| node[pos=0.25, above] {No} (F.north);
    \draw [arr] (E.south) |- (serviceJoin);
    \draw [arr] (F.south) |- (serviceJoin);
    \fill[GatewayArrow] (serviceJoin) circle (1.4pt);
    \draw [arr] (serviceJoin) -- (G.north);
    \draw [arr] (G.west) -| node[pos=0.25, above] {No} (H.north);
    \draw [arr] (G.east) -| node[pos=0.25, above] {Sí} (I.north);
    \draw [arr] (H.west) -- ++(-0.9,0) |- (C.west);
    \draw [arr] (I.east) -- ++(0.9,0) |- (B.east);
    \end{tikzpicture}
    }}
    \caption{Diagrama de flujo de continuidad operativa del servicio perimetral: publicación, retención local y recuperación automática mediante \textit{systemd} con almacenamiento y reenvío (elaboración propia).}
    \label{fig:flujo_rollback}
\end{figure}

\subsubsection{Secuencia operativa de recuperación}
A partir de este flujo se distinguen tres trayectorias de interés:

\begin{enumerate}
    \item \textbf{Flujo nominal:} cuando la publicación es correcta, el lote queda confirmado y se conserva una copia local como evidencia de respaldo.
    \item \textbf{Contingencia por indisponibilidad del destino:} cuando el intermediario \ac{MQTT} o el enlace no están disponibles, el lote se preserva en la cola \textit{JSONL}. Así, la pérdida deja de ser inmediata y el lote puede publicarse con posterioridad.
    \item \textbf{Contingencia por fallo del servicio:} si el proceso cae o se bloquea, \textit{systemd} reinicia la unidad y la cola local vuelve a cargarse, reduciendo el riesgo de perder telemetría ya normalizada.
\end{enumerate}

\subsubsection{Mitigación de errores humanos}
El esquema de continuidad no solo responde a fallos de red o del componente lógico; también amortigua errores humanos durante instalación, mantenimiento e investigación. La separación física mediante derivación controlada y aislamiento galvánico reduce la probabilidad de perturbar el bus original, mientras que la ejecución sin publicación externa, la lista de permitidos y la cola local disminuyen el impacto de una configuración incorrecta sobre el entorno productivo. De este modo, el diseño no elimina el error humano, pero limita su propagación hacia una pérdida silenciosa de evidencia o una publicación indebida.

\section{Motor de análisis sobre evidencia consolidada}
Una vez descritos los componentes que capturan, filtran, transportan y recuperan la telemetría en campo, el capítulo cierra con el motor de análisis, encargado de organizar, correlacionar e interpretar la evidencia consolidada durante las visitas técnicas.

Esta ubicación mantiene una separación nítida entre la pasarela, que opera durante la adquisición en campo, y el subsistema analítico, que trabaja después sobre datos ya consolidados. La solución no termina en la captura, pero ambos planos deben mantenerse separados: la pasarela actúa en campo; el motor ordena la evidencia y la convierte en insumos comparables para la validación.

\subsubsection{Etapas funcionales del motor de análisis}
La Figura~\ref{fig:fincadiag-flujo-balanceado} resume el flujo de análisis que recibe la evidencia una vez capturada. En este punto, cada dominio conserva su significado original hasta la fase de integración; así se evita confundir métricas de red con eventos seriales o inferir causalidad entre fenómenos observados en distintos momentos.

\begin{figure}[H]
\centering
\resizebox{0.993\linewidth}{!}{%
\begin{tikzpicture}[node distance=0.9cm and 1.2cm, transform shape, scale=0.70,
    startstop/.style={ellipse, draw=GatewayControlBorder, thick, fill=GatewayControlFill, minimum width=3.8cm, minimum height=1cm, align=center, font=\small\bfseries},
    process/.style={rectangle, rounded corners=3pt, draw=GatewayCoreBorder, thick, fill=GatewayCoreFill, text width=4.4cm, minimum height=1.15cm, align=center, font=\small},
    source/.style={rectangle, rounded corners=3pt, draw=GatewayExternalBorder, thick, fill=GatewayExternalFill, text width=4.4cm, minimum height=1.15cm, align=center, font=\small},
    warn/.style={rectangle, rounded corners=3pt, draw=GatewayDecisionBorder, thick, fill=GatewayDecisionFill, text width=4.4cm, minimum height=1.15cm, align=center, font=\small},
    decision/.style={diamond, aspect=2.0, draw=GatewayDecisionBorder, thick, fill=GatewayDecisionFill, text width=3.1cm, align=center, inner sep=2pt, font=\small},
    output/.style={trapezium, trapezium left angle=75, trapezium right angle=105, draw=GatewayExternalBorder, thick, fill=GatewayExternalFill, text width=4.0cm, minimum height=1.1cm, align=center, font=\small},
    arr/.style={-Latex, thick, GatewayArrow}
]

\node[startstop] (start) {Inicio};
\node[process, below=of start] (input) {Descubrimiento por visita y sesión};
\node[process, below=of input] (classify) {Clasificación de sesión\\ línea base, collar u ordeño completo};

\node[source, below left=1.4cm and 4.2cm of classify] (base) {Línea base de red\\ latencia, variación, ARP, rutas};
\node[source, below=1.6cm of classify] (serial) {Canal serial\\ demultiplexado, eventos y tandas};
\node[source, below right=1.4cm and 4.2cm of classify] (pcap) {Dominio PCAP y antena\\ puerto 6001, exposición y activos};

\node[process, below=2.0cm of serial] (engine) {Integración controlada\\ preservando significado};
\node[decision, below=1.1cm of engine] (corr) {¿Hay soporte\\ simultáneo?};
\node[warn, below left=1.4cm and 3.2cm of corr] (alerts) {Alertas\\ y contexto técnico};
\node[process, below right=1.4cm and 3.2cm of corr] (correlation) {Correlación temporal\\ y validación de campo};
\node[process, below=2.1cm of corr] (summary) {Consolidación\\ resumen por sesión, visita y fase};

\node[output, below left=1.2cm and 4.3cm of summary] (hourly) {Informe técnico\\ por sesión};
\node[output, below=1.2cm of summary] (visit) {Perfiles de Obj. 1\\ y preparación};
\node[output, below right=1.2cm and 4.3cm of summary] (global) {Resumen global\\ y panel};
\node[startstop, below=1.5cm of visit] (end) {Fin};

\draw[arr] (start) -- (input);
\draw[arr] (input) -- (classify);
\draw[arr] (classify) -- (base);
\draw[arr] (classify) -- (serial);
\draw[arr] (classify) -- (pcap);
\draw[arr] (base) |- (engine);
\draw[arr] (serial) -- (engine);
\draw[arr] (pcap) |- (engine);
\draw[arr] (engine) -- (corr);
\draw[arr] (corr.west) -| node[pos=0.2, above, font=\scriptsize] {No} (alerts.north);
\draw[arr] (corr.east) -| node[pos=0.2, above, font=\scriptsize] {Sí} (correlation.north);
\draw[arr] (alerts.south) |- (summary.west);
\draw[arr] (correlation.south) |- (summary.east);
\draw[arr] (summary) -- (hourly);
\draw[arr] (summary) -- (visit);
\draw[arr] (summary) -- (global);
\draw[arr] (hourly.south) |- (end.west);
\draw[arr] (visit) -- (end);
\draw[arr] (global.south) |- (end.east);

\end{tikzpicture}
}
\caption{Diagrama actual del instrumento FincaDiag Modular (elaboración propia).}
\label{fig:fincadiag-flujo-balanceado}
\end{figure}

Esta organización se puede leer como una secuencia de cinco etapas.

\begin{enumerate}
    \item \textbf{Descubrimiento y clasificación de sesión:} el motor recibe la carpeta de captura asociada a cada visita e identifica qué fuentes de evidencia están disponibles. Este primer paso define si la sesión aporta línea base, telemetría de collar o reconstrucción de ordeño completo.
    \item \textbf{Procesamiento por dominio:} cada fuente se trata primero en su propio plano. La línea base resume latencia, variación temporal y contexto de red; el serial se demultiplexa y se reorganiza en eventos y tandas; el dominio \ac{PCAP}/antena cuantifica exposición, firmas y cobertura del puerto 6001.
    \item \textbf{Integración controlada:} el motor normaliza marcas de tiempo, homogeneiza identificadores de sesión y preserva el origen de cada observación. La integración no borra el rastro de procedencia, porque la trazabilidad forma parte del resultado y no solo de la implementación técnica.
    \item \textbf{Correlación condicionada:} la correlación temporal solo se habilita cuando existe soporte simultáneo suficiente entre los dominios serial y de red. Cuando tal soporte no existe, el sistema entrega salidas descriptivas y alertas de contexto, pero evita forzar causalidad.
    \item \textbf{Consolidación y decisión de uso:} el motor produce resúmenes por sesión, visita y fase, junto con perfiles del Objetivo 1 y una lectura de preparación operativa de la pasarela. Esta última capa sirve para decidir qué evidencia ya respalda el diseño perimetral y cuál todavía requiere depuración.
\end{enumerate}

\subsubsection{Criterios de interpretación y trazabilidad}
El motor conserva la trazabilidad de cada sesión sin inferir comportamiento productivo cuando falta corroboración cruzada. Si solo existe evidencia serial o de red, entrega métricas descriptivas y explicita el límite de interpretación.

Como apoyo, puede enriquecer alertas con el inventario opcional \texttt{data/network/known\_hosts.json}, añadiendo nombres de dispositivos reconocidos sin modificar detecciones ni umbrales. Con ello se cierra el flujo: la pasarela preserva evidencia y el motor la organiza para evaluación posterior. La trazabilidad requisito--diseño--prueba queda distribuida entre capítulos: aquí se documenta el diseño y el Capítulo~\ref{chapter:chapter06} presenta su contraste de campo.
